\chapter{Surfaces}

Ce chapitre se concentre sur le fonctionnement de l'algorithme de Perlin \cite{PerlinNoise}.

\section{Bruit de Perlin}

\subsection{Notion de bruit}
Lors de l'application de l'algorithme de Perlin, le bruit est considéré comme une valeur pseudo-aléatoire générée par des coordonnées données en paramètres. 
La valeur du bruit pourra être interprétée à des étapes ultérieures lors de la construction d'un espace. 
% je n'aime pas cette formulation

\subsection{Notion de fonction de bruit}

Pour mettre en relation un ensemble de bruits, les fonctions de bruit font une interpolation entre deux valeurs, soit une moyenne. 
Selon la fonction choisie, l'interpolation sera plus au moins douce.
Dans le cas de l'algorithme de Perlin, nous utilisons une interpolation linéaire entre deux points.
Le paramètre \textit{alpha} permet de jouer 
\begin{verbatim}
float lerp(float x,float y,float alpha):
	return (1 - alpha)x + alpha * y
\end{verbatim}

\subsection{Fonctionnement de l'algorithme de Perlin} 

L'algorithme de Perlin peut se décomposer en plusieurs étapes : définir l'espace, générer un échantillon fini de valeurs aléatoires, faire des produits scalaires et ensuite interpoler. 



\subsection{Application}

\subsubsection{Mouvement brownien}


